\documentclass[12pt]{article}
\usepackage[utf8]{inputenc}
\usepackage[T1]{fontenc}
\usepackage{hyperref}
\title{Vumbu Language Resources for K-12 Curriculum Development}
\author{Compiled by Claude Sonnet 5 (claude-sonnet-5) for Andr\'as Kornai}
\date{}
\begin{document}
\maketitle
\section{Introduction}
Vumbu (ISO 639-3 vum; Glottolog vumb1238) is an endangered Bantu language of Gabon (Guthrie B.403, Sira group). A 2007 estimate gives about 5,000 speakers; Joshua Project reports about 11,000 people in the associated group, concentrated near Yetsou in the Moukalaba department, Ngouni\'e province. Dialect names reported are Bavoungou and Bawoumvou. Very little material exists, and no school use is known.
\section{Primary Educational Resources}
No primers, textbooks or school materials were found in search results.
\section{Linguistic and Reference Materials}
\begin{itemize}
\item Wikipedia overview of the Vumbu language: \url{https://en.wikipedia.org/wiki/Vumbu_language}
\item Glottolog entry: \url{https://glottolog.org/resource/languoid/id/vumb1238}
\end{itemize}
\section{Religious and Community Resources}
Joshua Project reports no Bible translation or Jesus Film in Vumbu and an unknown written status: \url{https://joshuaproject.net/people_groups/15804}
\section{Digital and Audio Resources}
Global Recordings Network lists the language (GRN 23102) but reports no Vumbu recordings; audio exists only for the related Bavoungou and Bawoumvou varieties: \url{https://globalrecordings.net/en/language/vum}
\section{Curriculum Development Context}
There is no written standard, no teaching materials and no known schooling in Vumbu. Any curriculum would start with documentation work and community consultation; speaker numbers are uncertain.
\end{document}